\documentclass[10pt]{scrartcl}
\usepackage[mdthm,colorsec,noauthor,noasy]{evan}
\usepackage{handout}
\newcommand{\aug}[2]{\left[\begin{array}{@{}#1@{}}#2\end{array}\right]}
\newcommand{\vv}[1]{\mathbf{#1}}
\DeclareMathOperator{\spn}{Span}\DeclareMathOperator{\Nul}{Nul}\DeclareMathOperator{\Col}{Col}
\title{Math 54 --- Week 5}
\subtitle{Solutions}
\author{}\date{}\ihead{}
\begin{document}
\maketitle

\begin{prob}\textbf{P1.}\par\idea{Swap to get a leading 1, use replacements only, count swaps.}
\begin{steps}
  \ii $R_1\leftrightarrow R_2$ (one swap): rows $(1,3,0,2)$, $(2,4,-2,6)$, $(-1,-1,4,1)$, $(3,8,1,5)$.
  \ii $R_2-2R_1$, $R_3+R_1$, $R_4-3R_1$: $(0,-2,-2,2)$, $(0,2,4,3)$, $(0,-1,1,-1)$.
  \ii $R_3+R_2$: $(0,0,2,5)$. $R_4-\tfrac12R_2$: $(0,0,2,-2)$.
  \ii $R_4-R_3$: $(0,0,0,-7)$. Diagonal: $1,-2,2,-7$, product $28$.
  \ii One swap: $\det=-28$.
\end{steps}
\end{prob}

\begin{prob}\textbf{P2.}
$\det(AB)=3(-2)=-6$. $\det(A^{-1})=\tfrac13$. $\det(A^T)=3$. $\det(2A)=2^3\cdot3=24$. $\det(B^TA^{-1}B)=(-2)\cdot\tfrac13\cdot(-2)=\tfrac43$.
\end{prob}

\begin{prob}\textbf{P3.}\par\idea{Invertible exactly when $\det\ne0$.}
Expand along row 1: $1(15-1)-2(6-0)+h(2-0)=14-12+2h=2(h+1)$. Invertible exactly when $h\ne-1$.
\end{prob}

\begin{prob}\textbf{P4.}
$\det A=5-6=-1$. $\det A_1(\vv b)=\det\begin{bmatrix}1&2\\4&1\end{bmatrix}=-7$, $\det A_2(\vv b)=\det\begin{bmatrix}5&1\\3&4\end{bmatrix}=17$. So $x_1=7$, $x_2=-17$. Check: $35-34=1$, $21-17=4$.
\end{prob}

\begin{prob}\textbf{P5.}\par\idea{Move one vertex to the origin; the two adjacent sides are the columns.}
Subtract $(-1,0)$: the vertices become $(0,0)$, $(1,5)$, $(2,-4)$, $(3,1)$, and $(1,5)+(2,-4)=(3,1)$, so the sides are $(1,5)$ and $(2,-4)$. Area $=|\det\begin{bmatrix}1&2\\5&-4\end{bmatrix}|=|-14|=14$. $\det\begin{bmatrix}2&1\\1&3\end{bmatrix}=5$, so the image has area $5\cdot14=70$.
\end{prob}

\begin{prob}\textbf{P6.}
\begin{steps}
  \ii (a) Yes: it is $\Nul\begin{bmatrix}1&1&0\end{bmatrix}$.
  \ii (b) No: $(1,0)$ and $(0,-1)$ are in it, but their sum $(1,-1)$ has $ab=-1<0$.
  \ii (c) No: the zero polynomial has value $0$ at $0$, not $1$.
  \ii (d) Yes: it is $\spn\{(1,1,0),(-2,0,3)\}$.
  \ii (e) Yes: $0^T=0$; $(A+B)^T=A^T+B^T=A+B$; $(cA)^T=cA^T=cA$.
\end{steps}
\end{prob}

\begin{prob}\textbf{P7.}
$\aug{rrr|r}{1&-3&-4&3\\-4&6&-2&3\\-3&7&6&-4}\to\aug{rrr|r}{1&0&5&-9/2\\0&1&3&-5/2\\0&0&0&0}$. Consistent, so $\vv b\in\Col A$. The coefficient part gives $x_1=-5x_3$, $x_2=-3x_3$: $\Nul A=\spn\{(-5,-3,1)\}$.
\end{prob}

\begin{prob}\textbf{P8.}
$A\vv u=(6-8+2+0,\ -6+10-7+0,\ 9-14+8+0)=(0,-3,3)\ne\vv0$: $\vv u\notin\Nul A$. Row reducing $[\,A\ \ \vv v\,]$ gives pivots in columns 1, 2, 4 only (not in the last column), so it is consistent: $\vv v\in\Col A$. $\vv u$ has 4 entries and $\Col A\subseteq\RR^3$, so $\vv u$ cannot be in $\Col A$.
\end{prob}

\end{document}
